\section{Proof that $\norm{\cdot}{1,k}$ is a semi-norm (Proposition~\ref{prop:seminorm})}
\label{app:seminorm}

\begin{proof}
    Absolute homonogeneity 
    $\norm{\lambda w}{1, k} = \absv{\lambda} \cdot \norm{w}{1,k}$
    is straightforward.
    We also have
    \begin{align*}
        \norm{w + w'}{1,k}
        &= \inf_{\substack{\bE \subset [E] \\ \card{\bE} = E - k + 1}}
            \sum_{e \in \bE} \absv{w_e + w_e'} 
        \leq \inf_{\substack{\bE \subset [E] \\ \card{\bE} = E - k + 1}} \set{
            \sum_{e \in \bE} \absv{w_e} 
            + \sum_{e \in \bE} \absv{w_e'} 
        } \\
        &\leq \inf_{\substack{\bE \subset [E] \\ \card{\bE} = E - k + 1}} \set{
            \sum_{e \in \bE} \absv{w_e}
        } + \inf_{\substack{\bE' \subset [E] \\ \card{\bE'} = E - k + 1}} \set{
            + \sum_{e \in \bE'} \absv{w_e'} 
        } 
        = \norm{w}{1,k} + \norm{w'}{1,k}.
    \end{align*}
    Note however that, while clearly nonnegative,
    $\norm{\cdot}{1,k}$ is not positive definite.
    Indeed, for $k \geq 2$, vectors ${\bf e}$ of the canonical basis 
    verify $\norm{\bf e}{1,k} = 0$.
\end{proof}


\section{Proof of \SimpleVeche{}'s resilience (Theorem~\ref{th:simpleveche})}
\label{app:simpleveche}

This section provides a proof of a generalization of Theorem~\ref{th:simpleveche},
which bounds the maximal impact of a followed user.
More precisely, we consider a slightly generalized setting with user volumes.
Define the weights as
\begin{align}
    \Weights(v, b) 
    \triangleq \max \set{0, \sum_{u \in [U]} v_u b_u / \norm{b_u}{1}}.
\end{align}

\begin{theorem}
\label{th:lipschitz_simpleveche}
    Let $v \in \setR_{\geq 0}^U$, $b \in \setR^{U \times E}$ and $k \in \setN$.
    If $\norm{W(v, b)}{1, k} > 0$, then
    \begin{align}
        \TotalVariationDistance \left(
            \SimpleVeche_{1:k} (v, b),
            \SimpleVeche_{1:k} (v', b),
        \right)
        \leq \norm{v - v'}{1} 
            \sum_{i=1}^k \norm{\Weights(v, b)}{1, i}^{-1}.
    \end{align}
\end{theorem}

Note that when 
$\norm{v - v'}{1} > \left( \sum \norm{\Weights(v, b)}{1, k}^{-1} \right)^{-1}$,
the bound becomes trivial as,
by definition, the total variational distance lies between $0$ and $1$.
This will hold in particular if 
$\norm{v - v'}{1} > \sum \norm{\Weights(v, b)}{1, k}$,
i.e. the shift in voting right allocations is larger 
than the worst-case weight allocation after $k-1$ recommendations.
It is also noteworthy that the obvious inequality
$\norm{w}{1,i} \leq \norm{w}{1,j}$ for $i < j$
implies that the right-hand side is itself upper-bounded by
$k \norm{v - v'}{1} \norm{\Weights (v, b)}{1,k}^{-1}$.

\begin{proof}[Proof of Theorem~\ref{th:simpleveche}]
Theorem~\ref{th:simpleveche} is easily derived 
from Theorem~\ref{th:lipschitz_simpleveche} as the special case
where $v \triangleq (1)_{u \in [U]} \in \setR_{\geq 0}^U$,
and $v'_{x} \triangleq v_x \indicator{x \neq u}$ for all $x \in [U]$,
i.e. $v'$ removes user $u$'s voting right.
\end{proof}


% \subsection{A few remarks}

% Before moving to the proof, let us make a couple of remarks.

% \paragraph{Weighing first recommendations more.}
% In practice, Alice may care more about the first recommendations that are made,
% rather than about the latter ones.
% While it does not do so formally, 
% our proof can be easily modified to account for this.
% More precisely, for $\tilde f, \tilde f' \in \cF_k^E$ 
% and $L \in \setR_{\geq 0}^k$, denote
% \begin{align}
%     \WeightedTotalVariationDistance (\tilde f, \tilde f' | L)
%     \triangleq \inf_{\text{Coupling} F, F' \text{ of } \tilde f, \tilde f'}
%         \sum_{i = 1}^k L_i \pb{F_i \neq F_i'}
% \end{align}


\paragraph{Relation with Lipschitz resilience.}
Theorem~\ref{th:lipschitz_simpleveche} looks similar 
to Lipschitz resilience~\cite{DBLP:conf/aistats/AllouahGHV24}.
Lipschitz resilience demands that 
the social mechanism $\setR_{\geq 0}^U \times \bbX^U \to \bbY$ 
be Lipschitz continuous as a function of users' voting rights
(using the $\ell_1$ norm as a metric for voting rights).
In our context, we would have $\bbX \triangleq \setR^E$ (user ballots)
and $\bbY \triangleq \Delta(\cF^E)$ with the $\TotalVariationDistance$ metric.
Theorem~\ref{th:lipschitz_simpleveche} does not guarantee Lipschitz resilience,
because our right-hand side is a function of $\norm{\Weights(v, b)}{1,i}$
for $i \in [1, k]$,
which depends on both $v$ and $b$.
This highlights two aspects of \SimpleVeche{}.

First, as already discussed in the main part of our paper,
the resilience of \SimpleVeche{} is weakened by negative ballot weights,
as preventing the recommendations of many entities
facilitate the arecommendations of other entities.
This could be mitigated 
by the prior identification of definitely safe entities to recommend.

Second, the dependence on $v$ simply highlights the fact 
that if Alice only follows a single user,
then her recommendations will be heavily influenced by this single user.
This dependency risk can be mitigated by encouraging Alice 
to multiply and diversify her list of followed accounts.
It could also be mitigated once again by a fallback to safe entities to recommend,
when her follow list is insufficiently long.


\subsection{Useful lemmas to prove Theorem~\ref{th:lipschitz_simpleveche}}

We first prove a lemma.

\begin{lemma}
    \label{lemma:weights_resilience}
    For all $v \in \setR_{\geq 0}^U$ and $b \in \setR^{U \times E}$,
    $\norm{\Weights(v, b) - \Weights(v', b)}{1} \leq \norm{v - v'}{1}$.
\end{lemma}

\begin{proof}
    Now note that, for all entity $e \in [E]$, we can bound
    \begin{align}
        \absv{w_e - w_e'}
        &= \absv{
            \max\set{0, \sum_{u \in [U]} v_u \tilde b_{ue}}
            - \max\set{0, \sum_{u \in [U]} v_u' \tilde b_{ue}}
        }.
    \end{align}
    Using the fact that $z \mapsto \max \set{0, z}$ 
    is clearly 1-Lipschitz continuous,
    we derive
    \begin{align}
        \absv{w_e - w_e'}
        &\leq \absv{
            \sum_{u \in [U]} v_u \tilde b_{ue}
            - \sum_{u \in [U]} v_u' \tilde b_{ue}
        }
        = \absv{\sum_{u \in [U]} (v_u - v_u') \tilde b_{ue}}
        \leq \sum_{u \in [U]} \absv{v_u - v_u'} \cdot \absv{\tilde b_{ue}}.
    \end{align}
    Adding up over $e \in [E]$ then yields
    \begin{align}
        \norm{w - w'}{1}
        &= \sum_{e \in [E]} \absv{w_e - w_e'}
        \leq \sum_{e \in [E]} \sum_{u \in [U]} \absv{v_u - v_u'} \cdot \absv{\tilde b_{ue}} \\
        &= \sum_{u \in [U]} \absv{v_u - v_u'} \sum_{e \in [E]} \absv{\tilde b_{ue}}
        = \sum_{u \in [U]}\absv{v_u - v_u'} \underbrace{\norm{\tilde b_u}{1}}_{=1}
        = \norm{v - v'}{1},
    \end{align}
    which yields the lemma.
\end{proof}

\begin{lemma}
    \label{lemma:one-sample}
    For $w \in \setR_{\geq 0}^E$ with $\norm{w}{1} > 0$,
    define $\tilde e(w) \in \Delta([E])$ the distribution of an element $e \in [E]$ selected,
    if it is drawn with probability $w_e / \norm{w}{1}$.
    Then, for any $w, w' \in \setR_{\geq 0}^E$ with $\norm{w}{1} > 0$ and $\norm{w'}{1} > 0$,
    \begin{align}
        \TotalVariationDistance(\tilde e(w), \tilde e(w')) \leq \frac{\norm{w - w'}{1}}{\norm{w}{1}}.
    \end{align}
\end{lemma}

\begin{proof}
    Given that $\norm{w}{1} > 0$ and $\norm{w'}{1} > 0$,
    the distributions of $\tilde e(w)$ and $\tilde e(w')$ 
    are discrete probabilities in $\Delta([E])$.
    It is well-known that the total variation distance for discrete probabilities is given by
    \begin{align}
        \TotalVariationDistance (\tilde e(w), \tilde e(w'))
        &= \frac{1}{2} \sum_{e^* \in [E]} \absv{
            \pb{e = e^* \st e \sim \tilde e(w)}
            - \pb{e = e^* \st e \sim \tilde e(w')}
        } \\
        &= \frac{1}{2} \sum_{e \in [E]} \absv{
            \frac{w_e}{\norm{w}{1}}
            - \frac{w_e'}{\norm{w'}{1}}
        }.
    \end{align}
    We now leverage the triangle inequality, yielding
    \begin{align}
        \TotalVariationDistance &(\tilde e(w), \tilde e(w'))
        = \frac{1}{2} \sum_{e \in [E]} \absv{
            \frac{w_e}{\norm{w}{1}} - \frac{w_e'}{\norm{w}{1}}
            + \frac{w_e'}{\norm{w}{1}} - \frac{w_e'}{\norm{w'}{1}}
        } \\
        &\leq \frac{1}{2} \sum_{e \in [E]} \frac{\absv{w_e - w_e'}}{\norm{w}{1}}
        + \frac{1}{2} \sum_{e \in [E]} \absv{
            \left( \frac{1}{\norm{w}{1}} - \frac{1}{\norm{w}{1}}\right) 
            w_e'
        } \\
        &\leq \frac{1}{2 \norm{w}{1}} \sum_{e \in [E]} \absv{w_e - w_e'}
        + \frac{\absv{\norm{w'}{1} - \norm{w}{1}}}{2 \norm{w}{1} \cdot \norm{w'}{1}} 
            \sum_{e \in [E]} \absv{w_e'} \\
        &= \frac{\norm{w - w'}{1} + \absv{\norm{w}{1} - \norm{w'}{1}}}{2 \norm{w}{1}}.
    \end{align}
    Again, by the triangle inequality, 
    we have $\norm{w - w'}{1} \geq \absv{\norm{w}{1} - \norm{w'}{1}}$.
    This then yields
    \begin{align}
        \TotalVariationDistance &(\tilde e(w), \tilde e(w'))
        \leq \frac{\absv{\norm{w}{1} - \norm{w'}{1}} + \norm{w - w'}{1}}{2 \norm{w}{1}} 
        \leq \frac{\norm{w - w'}{1}}{\norm{w}{1}},
    \end{align}
    which concludes the proof.
\end{proof}


\subsection{Total variation distance of joint distribution}

For any distribution $\tilde x$ on a finite set $\bbX$,
we define $\Support(\tilde x) \triangleq \set{x \in \bbX \st \pb{\tilde x = x} > 0}$
its support.

\begin{lemma}
    \label{lemma:joint-distribution}
    Consider two finite sets $\bbX_1$ and $\bbX_2$.
    Let $\tilde x, \tilde x' \in \Delta(\bbX_1 \times \bbX_2)$.
    Then 
    \begin{align}
        \TotalVariationDistance (\tilde x, \tilde x') 
        &\leq \TotalVariationDistance (\tilde x_1, \tilde x_1') 
        + \sup_{x_1 \in \Support(\tilde x_1) \cup \Support(\tilde x_1')} 
            \TotalVariationDistance (\tilde x_2 | x_1, \tilde x_2' | x_1).
    \end{align}
\end{lemma}

\begin{proof}
    By a slight abuse of notation, 
    we use interchangeably the distribution $\tilde x$
    and a random variable $\tilde x$ that follows this distribution.
    Denote alose $\bS_1 \triangleq \Support(\tilde x_1) \cup \Support(\tilde x_1')$.
    Our proof essentially boils down to the triangle inequality.
    More precisely, we have
    \begin{align}
        &\TotalVariationDistance (\tilde x, \tilde x') 
        = \frac{1}{2} \sum_{x \in \bbX_1 \times \bbX_2} \absv{
            \pb{\tilde x = x}
            - \pb{\tilde x' = x}
        } \\
        &= \frac{1}{2} \sum_{(x_1, x_2) \in \bbX_1 \times \bbX_2} \absv{
            \pb{\tilde x_1 = x_1 \wedge \tilde x_2 = x_2}
            - \pb{\tilde x_1' = x_1 \wedge \tilde x_2' = x_2}
        } \\
        &= \frac{1}{2} \sum_{(x_1, x_2) \in \bbX_1 \times \bbX_2} \absv{
            \pb{\tilde x_1 = x_1} \pb{\tilde x_2 = x_2 \st \tilde x_1 = x_1}
            - \pb{\tilde x_1' = x_1} \pb{\tilde x_2' = x_2 \st \tilde x_1' = x_1}
        } \\
        &\leq \frac{1}{2} \sum_{(x_1, x_2) \in \bbX_1 \times \bbX_2} \absv{
            \pb{\tilde x_1 = x_1} \pb{\tilde x_2 = x_2 \st \tilde x_1 = x_1}
            - \pb{\tilde x_1' = x_1} \pb{\tilde x_2 = x_2 \st \tilde x_1 = x_1}
        } \nonumber \\
        &\quad + \frac{1}{2} \sum_{(x_1, x_2) \in \bbX_1 \times \bbX_2} \absv{
            \pb{\tilde x_1' = x_1} \pb{\tilde x_2 = x_2 \st \tilde x_1 = x_1}
            - \pb{\tilde x_1' = x_1} \pb{\tilde x_2' = x_2 \st \tilde x_1' = x_1}
        } \\
        &\leq \frac{1}{2} \sum_{x_1 \in \bbX_1} 
        \absv{
            \pb{\tilde x_1 = x_1} - \pb{\tilde x_1' = x_1}
        } \underbrace{
            \sum_{x_2 \in \bbX_2} \pb{\tilde x_2 = x_2 \st \tilde x_1 = x_1}
        }_{= 1}
        \nonumber \\
        &\quad + \sum_{x_1 \in \bbX_1} 
        \pb{\tilde x_1' = x_1} \cdot \underbrace{
            \frac{1}{2} \sum_{x_2 \in \bbX_2} \absv{
                \pb{\tilde x_2 = x_2 \st \tilde x_1 = x_1}
                - \pb{\tilde x_2' = x_2 \st \tilde x_1' = x_1}
            }
        }_{= \TotalVariationDistance(\tilde x_2|x_1, \tilde x_2'|x_1)} \\
        &\leq \frac{1}{2} \sum_{x_1 \in \bbX_1} \absv{
            \pb{\tilde x_1 = x_1} - \pb{\tilde x_1' = x_1}
        } + \left( 
            \sup_{x_1 \in \bbX_1} 
            \TotalVariationDistance(\tilde x_2 | x_1, \tilde x_2' | x_1) 
        \right) \underbrace{\sum_{x_1 \in \bS_1} \pb{\tilde x_1' = x_1}}_{= 1} \\
        &= \TotalVariationDistance(\tilde x_1, \tilde x_1')
        + \sup_{x_1 \in \bS_1} 
            \TotalVariationDistance(\tilde x_2 | x_1, \tilde x_2' | x_1),
    \end{align}
    which is the lemma.
\end{proof}

% \begin{proof}
%     We construct a maximal coupling for two random variables $(X, X')$,
%     such that the marginal distribution of $X$ equals $\tilde x$,
%     and the marginal distribution of $X'$ equals $\tilde x'$.
%     It is well known that 
%     $\TotalVariationDistance (\tilde x, \tilde x') \leq \pb{X \neq X'}$.
%     Therefore
%     \begin{align}
%         \TotalVariationDistance (\tilde x, \tilde x') 
%         &\leq \pb{X_1 \neq X_1' \vee X_2 \neq X_2'} \\
%         &= \pb{X_1 \neq X_1' \vee \left( 
%             (X_2 \neq X_2') \wedge (X_1' = X_1) 
%         \right)} \\
%         &= \pb{X_1 \neq X_1'} + \pb{(X_2 \neq X_2') \wedge (X_1' = X_1)}.
%     \end{align}
% \end{proof}

\subsection{Proof of Theorem~\ref{th:lipschitz_simpleveche}}

\begin{proof}[Proof of Theorem~\ref{th:lipschitz_simpleveche}]
    Our proof goes by induction. 
    The case $k = 1$ is a direct result 
    of Lemma~\ref{lemma:weights_resilience} and 
    Lemma~\ref{lemma:one-sample}
    (using the fact that $\norm{\cdot}{1,1} = \norm{\cdot}{1}$).

    Now assume that 
    \begin{align}
        \TotalVariationDistance \left(
            \SimpleVeche_{1:k} (v, b),
            \SimpleVeche_{1:k} (v', b),
        \right)
        \leq \norm{v - v'}{1} \sum_{i=1}^k \norm{W(v, b)}{1, i}^{-1}.
    \end{align}
    We then note that $\SimpleVeche_{1:k+1}$ can be viewed as a distribution
    over $E^k \times E$,
    i.e. it first selects the first $k$ entities of the feed,
    and then the $(k+1)$-th entity.
    Denote $\cF^E_k \triangleq \set{f \in [E]^k \st \card{f([k])} = k}$
    the set of feeds of $k$ all-different entities
    and $d \triangleq \TotalVariationDistance \left(
        \SimpleVeche_{1:k+1} (v, b),
        \SimpleVeche_{1:k+1} (v', b),
    \right)$ the distance we want to upper-bound.
    By Lemma~\ref{lemma:joint-distribution}, we have
    \begin{align}
        d \leq \norm{v - v'}{1} \sum_{i=1}^k \norm{W(v, b)}{1, i}^{-1}.
            + \sup_{f_{1:k} \in \cF^E_k} \TotalVariationDistance\left(
            \SimpleVeche_{k+1} (v, b) | f_{1:k},
            \SimpleVeche_{k+1} (v', b) | f_{1:k},
        \right).
    \end{align}
    We now remark that $k+1$-th entity to be drawn 
    is drawn as described in Lemma~\ref{lemma:one-sample},
    but with the vectors $\Weights(v, b)_{|[E] - f_{1:k}([k])}$
    and $\Weights(v', b)_{|[E] - f_{1:k}([k])}$,
    i.e. with the weight vectors where the weights 
    of the previously drawn entities $e \in f_{1:k} ([k])$ are nullified.
    Therefore,
    \begin{align}
        d \leq \norm{v - v'}{1} \sum_{i=1}^k \norm{W(v, b)}{1, i}^{-1}.
        + \sup_{f_{1:k} \in \cF^E_k} \frac{
            \norm{
                \Weights(v, b)_{|[E] - f_{1:k}([k])}
                - \Weights(v', b)_{|[E] - f_{1:k}([k])}
            }{1}
        }{
            \norm{\Weights(v, b)_{|[E] - f_{1:k}([k])}}{1}
        }.
    \end{align}
    We now observe that canceling some coordinates 
    can only reduce the $\ell_1$ distance between two vectors, 
    thus
    $\norm{
        \Weights(v, b)_{|[E] - f_{1:k}([k])}
        - \Weights(v', b)_{|[E] - f_{1:k}([k])}
    }{1} \leq \norm{
        \Weights(v, b)
        - \Weights(v', b)
    }{1}$.
    By Lemma~\ref{lemma:weights_resilience},
    we know that this is at most $\norm{v - v'}{1}$.
    Hence
    \begin{align}
        d 
        &\leq \norm{v - v'}{1} \sum_{i=1}^k \norm{W(v, b)}{1, i}^{-1}.
        + \norm{v - v'}{1} \sup_{f_{1:k} \in \cF^E_k}
            \norm{\Weights(v, b)_{|[E] - f_{1:k}([k])}}{1}^{-1}
    \end{align}
    We now note that
    \begin{align}
        \sup_{f_{1:k} \in \cF^E_k}
            \norm{\Weights(v, b)_{|[E] - f_{1:k}([k])}}{1}^{-1}
        \leq \left(
            \inf_{\substack{\bF \subset [E] \\ \card{\bF} = k}}
            \norm{\Weights(v, b)_{|[E] - \bF}}{1}
        \right)^{-1}
        \leq \norm{\Weights(v, b)}{1, k+1}^{-1}.
    \end{align}
    Using additionally that we trivially have 
    $\norm{W(v, b)}{1, k} \geq \norm{W(v, b)}{1, k+1}$,
    we finally obtain
    \begin{align}
        d 
        \leq \norm{v - v'}{1} \sum_{i=1}^k \norm{W(v, b)}{1, i}^{-1}.
        + \frac{\norm{v - v'}{1}}{\norm{\Weights(v, b)}{1, k+1}}
        = \norm{v - v'}{1} \sum_{i=1}^{k+1} \norm{W(v, b)}{1, i}^{-1},
    \end{align}
    which concludes the proof.
\end{proof}


% \subsection{Paella}

% \begin{lemma}
%     \label{lemma:weights_resilience_removed_entities}
%     Let $v \in \setR_{\geq 0}^U$ and $b \in \setR^{U \times E}$.
%     Let $\bE \subset [E]$ and write 
%     $p_{\bE}: \setR^E \to \setR^E$ the orthogonal projection on 
%     $\set{ w \in \setR^E \st \forall e \notin \bE, w_e = 0}$
%     (i.e. $p_\bE$ sets to zero the coordinates that do note correspond to entities in $\bE$).
%     Then $\norm{p_\bE (\Weights(v, b)) - p_\bE (\Weights(v', b))}{1} \leq \norm{v - v'}{1}$.
% \end{lemma}

% \begin{proof}
%     This follows from the previous lemma and the fact that the projector is 1-Lipschitz continuous.
%     More precisely,
%     \begin{align}
%         \norm{p_\bE(w) - p_\bE(w')}{1}
%         = \sum_{e \in \bE} \absv{w_e - w'_e}
%         \leq \sum_{e \in [E]} \absv{w_e - w'_e}
%         = \norm{w - w'}{1}.
%     \end{align}
%     Plugging $w \triangleq \Weights(v, b)$ and $w' \triangleq \Weights(v', b)$ yields the lemma.
% \end{proof}